% SCI-RTC v0.1/r0.3 shared normative core -- equations.
\section{Normative equations and identities}
\label{sec:core-equations}

The equations in this section are the sole displayed mathematical authority
of both views.  They describe an implementation-independent, state-conditioned
operator.  They neither select numerical policy nor assert implementation
conformity.

\subsection{Role domain and donor transfer}

For an exact detector occurrence $i$ under the admitted multiplicative
raw-to-calibrated convention,
\begin{equation}
 z_i=\phi_i x_i,
 \qquad \phi_i=\texttt{flxscale}_i .
 \tag{SCI-RTC-EQ-001}\label{eq:raw-cal}
\end{equation}
Equation~\eqref{eq:raw-cal} does not assign a role to legacy
\texttt{responsivity}.  When $q$ is a raw donor and $d$ is a raw target, the
unique multiplicative transfer that represents donor-calibrated amplitude in
the target raw convention is
\begin{equation}
 x^{R}_{d\leftarrow q}
   =a_{d\leftarrow q}x_q,
 \qquad
 a_{d\leftarrow q}=\frac{\phi_q}{\phi_d},
 \qquad
 \phi_d x^{R}_{d\leftarrow q}=\phi_qx_q .
 \tag{SCI-RTC-EQ-002}\label{eq:raw-donor}
\end{equation}
This equation is defined only for a compatible \texttt{flxscale} pair as
defined above.  Otherwise this transfer is unavailable.  An
additive calibration convention, a different unit/reference domain, or an
alternative donor authority requires its own complete affine equivalence and
is not silently folded into equation~\eqref{eq:raw-donor}.

Let the role-domain operator be
\begin{equation}
 C_{\mathcal R}=\begin{cases}
 I, & \delta_{\mathcal R}=0\quad\text{(admitted raw role)},\\
 C_\kappa, & \delta_{\mathcal R}=1\quad\text{(CAL-authorized role)}.
 \end{cases}
 \tag{SCI-RTC-EQ-003}\label{eq:role-cal}
\end{equation}
For a CAL-authorized role, calibration precedes cross-detector replacement.
For a raw role, $C_{\mathcal R}=I$ and an admitted raw donor transfer may use
equation~\eqref{eq:raw-donor}.  Reversing calibration and replacement is
equivalent only if, for the complete admitted affine operators on the exact
support,
\begin{equation}
 B^{(z)}_{\omega} C_\kappa
   =C_\kappa B^{(x)}_{\omega},
 \quad
 \mathbf b^{(z)}_{\omega}
   =C_\kappa\mathbf b^{(x)}_{\omega}
     +\mathbf c_\kappa-B^{(z)}_{\omega}\mathbf c_\kappa,
 \tag{SCI-RTC-EQ-004}\label{eq:order-equivalence}
\end{equation}
with unit, factor direction, additive offset, response, uncertainty, identity,
validity, and domain equivalence all proved.  Equality of numerical values in
one fixture is insufficient.

\subsection{Realized conditioning operator}

Conditioned on complete realized state $\Omega$, the ordered role-specific
operator is
\begin{equation}
 \boxed{\;
 \mathbf y_{\mathcal R}=D_{M,0}
 F_{K,\theta_K}\cdots F_{1,\theta_1}
 B_{\omega,\mathcal R}
 \left[C_{\mathcal R}
  (A_\alpha\mathbf x^{\rm acq}+\mathbf a_\alpha)
  +\mathbf c_{\mathcal R}\right]
 +\mathbf b_{\Omega}\; .}
 \tag{SCI-RTC-EQ-005}\label{eq:complete-operator}
\end{equation}
Chronological order is read from right to left.  The affine terms include
only declared calibration, replacement, filter-state, boundary, or other
authorized contributions.  Masks, coefficients, finite handling, context,
exceptions, and state are components of $\Omega$, not annotations outside
the operator.  Conditional on them,
\begin{equation}
 \mathbf y_{\mathcal R}=L_\Omega\mathbf x^{\rm acq}
  +\mathbf b'_\Omega,
 \qquad
 L_\Omega=D_{M,0}F_K\cdots F_1B_{\omega,\mathcal R}
 C_{\mathcal R}A_\alpha .
 \tag{SCI-RTC-EQ-006}\label{eq:affine-operator}
\end{equation}
This is a conditional affine identity, not an unconditional linear,
stationary, unbiased, or Gaussian claim.

A realized replacement at target cell $(d,j)$ may be written
\begin{equation}
 u^R_{dj}=(1-p_{dj})u_{dj}
 +p_{dj}\left[
 \sum_{(q,\ell)\in\mathcal Q_{dj}}
 w_{dj,q\ell}\,a_{d\leftarrow q,j\ell}\,u_{q\ell}
 +\beta_{dj}\right],
 \quad
 \sum_{q,\ell}w_{dj,q\ell}=1
 \ \text{when normalization is claimed}.
 \tag{SCI-RTC-EQ-007}\label{eq:replacement}
\end{equation}
The scale $a$ is domain-qualified: equation~\eqref{eq:raw-donor} is the
authorized raw \texttt{flxscale} route, while a calibrated-domain transfer or
any separately authorized route carries its own identity and validity.

For a time-invariant FIR stage,
\begin{equation}
 v_{d,j}=\sum_{\ell\in\mathcal L_k}h^{(k)}_{d,\ell}u_{d,j-\ell},
 \qquad
 H^{(k)}_d(e^{i\omega})=
 \sum_{\ell\in\mathcal L_k}h^{(k)}_{d,\ell}e^{-i\omega\ell},
 \qquad H^{(k)}_d(1)=\sum_\ell h^{(k)}_{d,\ell}.
 \tag{SCI-RTC-EQ-008}\label{eq:fir}
\end{equation}
For a causal IIR or notch stage, a state-complete form is
\begin{equation}
 \mathbf s_{j+1}=A_k\mathbf s_j+B_ku_j,
 \qquad v_j=C_k\mathbf s_j+D_ku_j,
 \qquad \mathbf s_{j_0}=\mathbf s_{0,k}.
 \tag{SCI-RTC-EQ-009}\label{eq:iir}
\end{equation}
An exactly equivalent recurrence is admissible only with its coefficient
convention, precision, state, and boundary identity.  A causal-stability claim
requires every realized pole strictly inside the unit circle.

One possible selected mask operator is normalized masked FIR convolution,
\begin{equation}
 v_j=
 \frac{\sum_\ell h_\ell m_{j-\ell}u_{j-\ell}}
      {\sum_\ell h_\ell m_{j-\ell}},
 \qquad
 \left|\sum_\ell h_\ell m_{j-\ell}\right|>\epsilon_m ,
 \tag{SCI-RTC-EQ-010}\label{eq:masked-fir}
\end{equation}
but equation~\eqref{eq:masked-fir} is not a selected default.  Segmenting,
zero filling, restoring a mask after filtering, holding state, and
renormalizing are distinct operators.

\subsection{Phase-zero sampling and response}

For a coherent pre-sampling segment of length $N$, v0.1 selects only
\begin{equation}
 y_{d,n}=u_{d,Mn},
 \qquad 0\le Mn<N,
 \qquad f_{\rm out}=\frac{f_{\rm in}}{M},
 \qquad
 N_y=\begin{cases}0,&N=0,\\1+\left\lfloor\dfrac{N-1}{M}\right\rfloor,&N>0.
 \end{cases}
 \tag{SCI-RTC-EQ-011}\label{eq:phase-zero}
\end{equation}
The representative assigned-grid time is the selected point time
$t^{\rm out}_n=t_{Mn}$; the scientific support is nevertheless the complete
transitive support $\mathcal S_{dn}$ of all preceding stages.  A selected
edge/output policy may reject candidate points, but shall not relabel a
different phase or silently change equation~\eqref{eq:phase-zero}.

For output $(d,n)$ and acquired input $(q,j)$, the exact local response is
\begin{equation}
 K_{dn,qj}=\frac{\partial y_{dn}}
                 {\partial x^{\rm acq}_{qj}}
 \quad\text{at fixed realized state }\Omega .
 \tag{SCI-RTC-EQ-012}\label{eq:local-kernel}
\end{equation}
At a selection boundary the derivative may be discontinuous or undefined;
that state shall be represented rather than replaced by an LTI claim.

In the fixed, detector-separable, translation-invariant, single-rate interior
case only,
\begin{equation}
 H^{\rm comb}_d(e^{i\omega})=
 \prod_{k=1}^{K}H^{(k)}_d(e^{i\omega}),
 \qquad
 \tau_g(\omega)=-\Delta t\,\frac{d}{d\omega}
 \operatorname{unwrap}\!\left[\arg H^{\rm comb}_d(e^{i\omega})\right].
 \tag{SCI-RTC-EQ-013}\label{eq:lti-response}
\end{equation}
Amplitude is $|H|$, power response is $|H|^2$, and phase is $\arg H$ under
the recorded Fourier convention.  Group delay is unavailable at response
zeros unless a bounded convention is supplied.

For an LTI prefilter followed by phase-zero selection, the exact multirate
identity is
\begin{equation}
 Y(e^{i\Omega_o})=\frac{1}{M}\sum_{r=0}^{M-1}
 H\!\left(e^{i(\Omega_o+2\pi r)/M}\right)
 X\!\left(e^{i(\Omega_o+2\pi r)/M}\right),
 \tag{SCI-RTC-EQ-014}\label{eq:alias}
\end{equation}
under the declared Fourier/sign convention.  Every folded image contributes;
point selection itself is not anti-alias filtering.

For an FIR output, direct support starts with
\begin{equation}
 \mathcal S^{(1)}_{dn}=\{Mn-\ell:\ell\in\mathcal L\},
 \qquad
 \mathcal S_{dn}=\operatorname{closure}_{A,B,F,\,\mathrm{state}}
                 \!\left(\mathcal S^{(1)}_{dn}\right).
 \tag{SCI-RTC-EQ-015}\label{eq:support}
\end{equation}
IIR support extends to the declared state/reset or context boundary and may
be mathematically infinite.  Any finite guard is an approximation with a
stated residual-tail criterion.

\subsection{Moments, influence, and learned-plan selection}

For admitted input mean $\boldsymbol\mu_x$ and full covariance $\Sigma_x$,
conditional on fixed $\Omega$ independent of the modeled random noise,
\begin{align}
 \E[\mathbf y\mid\Omega]
   &=L_\Omega\boldsymbol\mu_x+\mathbf b'_\Omega,
 &\tag{SCI-RTC-EQ-016a}\label{eq:conditional-mean}\\
 \Sigma^{\rm stat}_{y\mid\Omega}
   &=L_\Omega\Sigma_xL_\Omega^{\mathsf T}.
 &\tag{SCI-RTC-EQ-016b}\label{eq:conditional-cov}
\end{align}
If selectors depend on the random data, unconditional covariance obeys
\begin{equation}
 \Cov(\mathbf y)=
 \E_\Omega[\Cov(\mathbf y\mid\Omega)]
 +\Cov_\Omega(\E[\mathbf y\mid\Omega]).
 \tag{SCI-RTC-EQ-017}\label{eq:total-covariance}
\end{equation}
The second term is not zero by default.  With named nuisance parameters
$\boldsymbol\eta$ and an adequate local linearization,
\begin{equation}
 \Sigma_y^{\rm sys}\approx
 J_\eta\Sigma_\eta J_\eta^{\mathsf T},
 \qquad
 J_\eta=\left.
 \frac{\partial\E[\mathbf y\mid\Omega,\boldsymbol\eta]}
      {\partial\boldsymbol\eta}\right|_{\boldsymbol\eta_0},
 \tag{SCI-RTC-EQ-018}\label{eq:systematic-cov}
\end{equation}
and a total-covariance claim requires every applicable component and cross
term:
\begin{equation}
 \Sigma_y^{\rm total}=
 \Sigma_y^{\rm stat}+\Sigma_y^{\rm sys}
 +\Sigma_y^{\rm selection}+\Sigma_y^{\rm model}
 +\Sigma_y^{\rm cross}.
 \tag{SCI-RTC-EQ-019}\label{eq:total-uncertainty}
\end{equation}
Unknown components are unavailable, not zero.

Define cause-preserving influence and the v0.1 ineligibility rule by
\begin{align}
 \mathcal I_{dn}
   &=\operatorname{closure}_{K,\,\mathrm{generative\ links}}
     \{\text{typed source causes on }\mathcal S_{dn}\},
 &\tag{SCI-RTC-EQ-020a}\label{eq:influence}\\
 E^{\rm RTC}_{dn}
   &=V^{\rm det}_d\land V^{\rm num}_{dn}\land
     V^{\rm support}_{dn}\land V^{\rm response}_{dn}
     \land\neg\mathbb I[\mathrm{ALIGN\_synth}\in\mathcal I_{dn}]
     \land\neg\mathbb I[\mathrm{RTC\_replace}\in\mathcal I_{dn}].
 &\tag{SCI-RTC-EQ-020b}\label{eq:eligibility}
\end{align}
Consumer-specific coordinate or policy predicates may further restrict
eligibility; they may not remove either transitive-influence exclusion.

For learned mode, let $\mathcal M$ be the admitted finite integer candidate
set and $\mathcal F(M)$ the conjunction of all owner-resolved astronomical
transfer, exact phase-zero alias, sampling, realized-filter realizability,
and downstream-compatibility constraints, evaluated with the smallest
admitted beam and maximum valid in-scan speed.  Then
\begin{equation}
 \mathcal M_{\rm safe}=\{M\in\mathcal M:\mathcal F(M)=\mathrm{true}\},
 \qquad M_\star=\max\mathcal M_{\rm safe}.
 \tag{SCI-RTC-EQ-021}\label{eq:learned-plan}
\end{equation}
Equation~\eqref{eq:learned-plan} is defined only after every numerical
tolerance, support rule, fallback, and compatibility predicate named in the
owner ledger is resolved.  Percentile speed is diagnostic, not safety
authority.

Finally, the atomic product is
\begin{equation}
 \mathcal O_{\rm RTC,\mathcal R}=
 (\mathbf y,\mathcal J,\mathcal K,\mathcal S,\mathcal I,
  \mathcal F,\mathcal V,\mathcal U,\mathcal P,\mathcal D,\mathcal C),
 \tag{SCI-RTC-EQ-022}\label{eq:bundle}
\end{equation}
where $\mathcal J$ is ordered identity, $\mathcal K$ complete response status,
$\mathcal S$ support, $\mathcal I$ influence, $\mathcal F$ typed flags/causes,
$\mathcal V$ separated validity and eligibility inputs, $\mathcal U$
uncertainty/covariance availability, $\mathcal P$ one-way state/provenance,
$\mathcal D$ diagnostics, and $\mathcal C$ atomic completion or required-output
failure state.

\subsection{Scientific design and learn--resolve--apply identities}

For a locally Gaussian effective beam with angular covariance $\Sigma$ and
angular scan-velocity vector $\mathbf v$, the projected source-crossing time
scale is
\begin{equation}
 \sigma_t=\left(\mathbf v^{\mathsf T}\Sigma^{-1}\mathbf v\right)^{-1/2}.
 \tag{SCI-RTC-EQ-023}\label{eq:crossing-time}
\end{equation}
Equation~\eqref{eq:crossing-time} is a local Gaussian diagnostic.  Scientific
admission shall use the exact source-times-realized-response model over the
declared beam, scan, source-class, and cadence domain.

A generic realized notch may be represented, where applicable, by
\begin{equation}
 H_{\rm notch}(z)=g\,
 \frac{\prod_{a=1}^{N_z}(1-z_a z^{-1})}
      {\prod_{b=1}^{N_p}(1-p_b z^{-1})},
 \qquad |p_b|<1\ \text{for a causal-stability claim}.
 \tag{SCI-RTC-EQ-024}\label{eq:notch-response}
\end{equation}
This form imposes no universal biquad.  Exact realized coefficients, sample
rate, frequency convention, phase/direction, state, edges, and precision are
part of the response.

For an admitted FIR design family $H_N$, a selected order shall follow a
declared constrained problem of the form
\begin{equation}
 N_\star=\min_N\left\{N:\begin{array}{l}
 \sup_{f\in\mathcal P}|H_N(f)-H_{\rm target}(f)|\le\epsilon_{\mathcal P},\\
 \sup_{f\in\mathcal S}|H_N(f)|\le\epsilon_{\mathcal S},\\
 \text{phase/delay, alias, edge/guard, precision, and realizability predicates pass}
 \end{array}\right\}.
 \tag{SCI-RTC-EQ-025}\label{eq:fir-order-selection}
\end{equation}
The sets, metrics, tolerances, equality conventions, and family are
owner-resolved; an approximate order formula is not universal authority.

For a causal linear-phase FIR with $N$ taps at rate $f_{\rm samp}$, its
nominal interior group delay is
\begin{equation}
 \tau_g=\frac{N-1}{2f_{\rm samp}}.
 \tag{SCI-RTC-EQ-026}\label{eq:fir-group-delay}
\end{equation}
Where a temporal response has effective delay $\tau$ during a locally uniform
scan, the corresponding scan-direction angular displacement is
\begin{equation}
 \Delta\boldsymbol\theta=\mathbf v\,\tau .
 \tag{SCI-RTC-EQ-027}\label{eq:scan-displacement}
\end{equation}

For every learned RTC operation, the lifecycle is
\begin{equation}
 \Lambda=\operatorname{Learn}(\mathcal L;\Theta^{\rm eff}),\qquad
 \Pi=\operatorname{Resolve}(\Lambda,\Theta^{\rm policy}),\qquad
 \mathbf y=\mathcal A_{\Pi}(\mathbf x),\qquad
 \frac{\partial\Pi}{\partial j}=0\ \text{during apply}.
 \tag{SCI-RTC-EQ-028}\label{eq:learn-resolve-apply}
\end{equation}
The final equality denotes plan immutability over apply slots; it is not an
assumption that the realized operator is stationary at masks, edges, state
boundaries, or other declared exceptions.

For bounded iterative notch-plan refinement with original admitted input
$x^{(0)}$, initial fixed or identity plan $\Pi_0$, and finite $K$,
\begin{equation}
 \begin{aligned}
 \Lambda_k&=\operatorname{Learn}(\mathcal L_k,
             x^{\rm eval,(k-1)},\Pi_{k-1}),\\
 \Pi_k&=\operatorname{Resolve}(\Lambda_k,\Pi_{k-1},
             \Theta^{\rm policy}),\\
 x^{\rm eval,(k)}&=\mathcal A_{\Pi_k}(x^{(0)}),\qquad
 x^{\rm final}=\mathcal A_{\Pi_K}(x^{(0)}),\\
 \frac{\partial\Pi_k}{\partial j}&=0
             \quad\text{during the apply slots of cycle }k,
 \qquad 1\le k\le K<\infty .
 \end{aligned}
 \tag{SCI-RTC-EQ-029}\label{eq:iterative-refinement}
\end{equation}
Equation~\eqref{eq:iterative-refinement} states the default original-input
replay rule.  A sequential cascade from $x^{\rm eval,(k-1)}$ requires separate
explicit operator authority and complete propagation of every intermediate
parent, response, state, phase, edge, support, covariance, and uncertainty.
